\maketitle
\thispagestyle{empty}

% =============================================================================
% ABSTRACT
% =============================================================================
\begin{abstract}
\noindent
This specification defines the full implementation of the \textbf{Simula} reasoning-driven
synthetic training data generation framework for SAP HANA Cloud. The system addresses a
critical gap identified in the existing training pipeline: \emph{there is no direct pipe
to HANA to generate training datasets for model training}. Three orthogonal components
are specified in full:

\begin{itemize}[nosep]
    \item a \textbf{HANA Schema Extraction Pipeline} that connects to SAP HANA Cloud via
          the AI Core PAL MCP service, queries \texttt{SYS.TABLE\_COLUMNS} for live metadata,
          and populates a schema registry;
    \item a \textbf{Taxonomy Generation Engine} that implements Algorithm~1 from the Simula
          research paper---factor identification, breadth-first expansion with Best-of-N
          proposals, and generator-critic refinement;
    \item a \textbf{Training Data Generation Engine} that implements Algorithm~2---taxonomic
          mix sampling, meta-prompt generation, complexification, and double-critic filtering.
\end{itemize}

\noindent
All three components consume a single vLLM TurboQuant backend via the OpenAI-compatible API
and output JSONL training data suitable for LoRA fine-tuning. The specification targets
three application domains: Arabic document processing (invoice extraction), Trial Balance
review (variance analysis and LLM commentaries), and ESG analytics (financed emissions
classification).

\medskip
\noindent\textbf{Reading order.} Chapter~\ref{chap:overview} situates the work;
Chapter~\ref{chap:schema} defines the data structures; Chapter~\ref{chap:extraction}
covers the HANA extraction pipeline; Chapter~\ref{chap:taxonomy} covers the taxonomy
engine; Chapter~\ref{chap:generator} covers the data generator; Chapter~\ref{chap:config}
specifies environment and CLI; Chapter~\ref{chap:evaluation} covers the evaluation
framework; Chapter~\ref{chap:references} provides references and traceability.

\medskip
\noindent\textbf{Conventions.} File paths are typeset in \texttt{monospace} and are
relative to the repository root \texttt{/Users/user/Documents/sap-oss}. Enum values
in running text appear in \textsc{small caps}. Python code listings include line numbers
for cross-reference.
\end{abstract}

\tableofcontents
\listoftables
\listoffigures
\newpage

% =============================================================================
% EXECUTIVE SUMMARY
% =============================================================================
\section*{Executive Summary}
\addcontentsline{toc}{section}{Executive Summary}

\subsection*{Business Problem}

The SAP OSS Studio AI platform requires high-quality training data to fine-tune 
Text-to-SQL models for domain-specific applications (Trial Balance, Arabic invoices, 
ESG analytics). The current pipeline uses \textbf{static Excel files} as schema sources,
creating three critical issues:

\begin{enumerate}[nosep]
    \item \textbf{Schema drift:} Training data becomes stale when HANA schemas evolve.
    \item \textbf{Coverage gaps:} Manual templates miss edge cases and rare query patterns.
    \item \textbf{Quality variance:} No systematic quality assurance for generated examples.
\end{enumerate}

\subsection*{Solution: Simula Framework}

This specification implements the \textbf{Simula} framework from Davidson et al. (2026),
a research-validated methodology for reasoning-driven synthetic data generation. The
framework addresses all three issues through:

\begin{itemize}[nosep]
    \item \textbf{HANA-Direct Extraction} — Live schema queries via AI Core PAL MCP service.
    \item \textbf{Taxonomic Coverage} — Hierarchical factor trees ensure systematic diversity.
    \item \textbf{Agentic Quality Control} — Double-critic filtering with Elo-calibrated complexity.
\end{itemize}

\subsection*{Key Metrics and Targets}

\begin{table}[H]
\centering
\begin{tabular}{llll}
\toprule
\textbf{Metric} & \textbf{Paper Target} & \textbf{Implementation} & \textbf{Source} \\
\midrule
Global diversity & 0.48--0.54 cosine & Target: 0.52 & Figure 4 \\
Local diversity & 0.10--0.30 cosine & Target: 0.20 & Figure 4 \\
Complexity ratio & 50\% (c=0.5) & Adaptive 0.2--0.5 & Figure 7 \\
Critic acceptance & p(y) $>$ 0.8 & Threshold: 0.5--0.7 & Figure 3 \\
Taxonomy completeness & 74\% (Simula) & Min: 70\% & Table 2 \\
Taxonomy soundness & 97\% (Simula) & Min: 90\% & Table 2 \\
Saturation ratio & 83\% & 83\% & Section 4.3 \\
\bottomrule
\end{tabular}
\end{table}

\subsection*{Expected Outcomes}

\begin{enumerate}[nosep]
    \item \textbf{10$\times$ training data} with consistent quality (100k+ examples).
    \item \textbf{Full taxonomic coverage} — no systematic gaps in query patterns.
    \item \textbf{Calibrated complexity} — downstream model benefits from high-complexity data.
    \item \textbf{Audit trail} — every example traceable to taxonomy nodes and meta-prompts.
\end{enumerate}

\newpage
% =============================================================================
% CHAPTER 1: OVERVIEW AND SCOPE
% =============================================================================